% Generated draft: requires mathematical and source review.
\documentclass[11pt]{article}
\usepackage[margin=25mm]{geometry}
\usepackage{fontspec}
\setmainfont{Latin Modern Roman}
\usepackage{amsmath,amssymb,mathtools,unicode-math}
\setmathfont{Latin Modern Math}
\usepackage{enumitem,needspace,xurl,hyperref}
\setlength{\parindent}{0pt}
\setlength{\parskip}{5pt}
\setlength{\emergencystretch}{3em}
\newcommand{\N}{\mathbb{N}}
\newcommand{\Z}{\mathbb{Z}}
\newcommand{\Q}{\mathbb{Q}}
\newcommand{\R}{\mathbb{R}}
\newcommand{\C}{\mathbb{C}}
\newcommand{\F}{\mathbb{F}}
\newcommand{\A}{\mathbb{A}}
\newcommand{\PP}{\mathbb{P}}
\newcommand{\E}{\mathbb{E}}
\newcommand{\eps}{\varepsilon}
\newcommand{\dd}{\,\mathrm{d}}
\begin{document}
\section*{UnsolvedMath ID: 2530}
\textbf{Source problem code:} catalogue entry
\subsection{Definitions and mathematical statement}
Let $k$ be an algebraically closed field and let $M$ be an algebraic group over $k$. The source explicitly calls $M$ simple in the sense that it has no proper nontrivial connected normal algebraic subgroup. Write $M(k)$ for its group of $k$-points and
\[
 Z(M(k))=\{z\in M(k):zm=mz\text{ for every }m\in M(k)\}.
\]
In the group-action formulation, $M/Z(M)$ is read as the corresponding central quotient of the point group. The source does not separately state connectedness or noncommutativity of $M$; its given simplicity clause is retained rather than silently strengthened.

Let $S$ be a simple group of finite Morley rank. For a group in a first-order structure, finite Morley rank means its underlying definable set has finite rank in the model-theoretic sense: nonempty sets have rank at least zero, and rank increases when infinitely many pairwise disjoint definable subsets of the preceding rank exist, with the usual limit-ordinal clause and parameters allowed. The rank and definability depend on the ambient structure. Ordinary abstract simplicity means no proper nontrivial normal subgroup; the source does not say whether a definable-simplicity convention is intended for this target group.

To speak of a definable action, $M(k)$, $S$ and the action map must be presented in a common first-order setting. For a group action $\rho:M(k)\times S\to S$, the identities are
\[
 \rho(1,s)=s,\qquad \rho(mm',s)=\rho(m,\rho(m',s)).
\]
Faithfulness means that $\rho(m,s)=s$ for every $s\in S$ implies $m=1$. If action by group automorphisms is intended, one must additionally require $\rho(m,ss')=\rho(m,s)\rho(m,s')$ for all $m,s,s'$. The wording ``on a group'' suggests that interpretation, but it is not explicitly specified.

\textbf{Imported rigidity assertion.} Prove that a definable, faithful, irreducible action of $M$ on such an $S$ can occur only when
\[
 S\cong M(k)/Z(M(k)).
\]
Equivalently, no such action should exist for a simple finite-Morley-rank target outside that isomorphism class. This identifies the source's exception without asserting existence or uniqueness of an exceptional action. In particular, the natural conjugation action on the central quotient factors through that quotient and need not be faithful as an action of $M(k)$ when the center is nontrivial.

The source does not define irreducibility. For an automorphism action, one possible convention excludes proper nontrivial definable invariant subgroups; another restricts attention to connected definable invariant subgroups. A set-action notion is different again. None is selected here as the uniquely intended hypothesis. Likewise, the source does not say whether the concluding isomorphism must be definable or merely abstract, or fix the ambient structure and parameter convention. These omissions prevent a complete theorem-level specification without review.

The literature triage dated 2026-08-17 reports no verified proof of the requested nonexistence assertion and labels it open. Its observation that $M/Z(M)$ is the expected exception is a description of the statement, not a proved classification result. The registry also labels the problem open. The discussion/literature field names A. V. Borovik, and the source identifies Kourovka Notebook, New Problems, 21st issue (2026), Problem 21.21. No independent source consultation or current-status assessment is claimed.
\subsection{Short English statement}
The proposed rigidity statement says that a simple algebraic group cannot act definably, faithfully and irreducibly on a simple group of finite Morley rank except for its own central quotient. The source leaves the ambient definability setting, automorphism-action convention, target simplicity and irreducibility insufficiently specified, so those hypotheses need clarification before a precise theorem can be adopted.
\subsection{Sources}
\begin{itemize}
\item UnsolvedMath local registry, data/unsolvedmath/problems.json, permanent ID 2530, source code catalogue entry, ``Kourovka Notebook Problem 21.21,'' including the explicit algebraic simplicity clause and literature triage dated 2026-08-17.
\item \emph{Kourovka Notebook}, New Problems, 21st issue (2026), Problem 21.21, supplied site \url{https://www.kourovka-notebook.org/}. The supplied review identifies this as the maintained source of the proof problem; no direct problem page, theorem-level reference or exact surrounding definitions are supplied.
\item A. V. Borovik, as named in the supplied discussion/literature field. No separate title or bibliographic entry is given, and no authorship details are invented.
\item Ulam AI and the UnsolvedMath Contributors, \emph{UnsolvedMath: A Curated Collection of Open Mathematics Problems}. Dataset curation and original metadata are supplied under CC BY 4.0; referenced source material retains its own terms. No external source was independently consulted for this draft.
\end{itemize}
\end{document}
